\documentclass[11pt]{article}
\usepackage[utf8]{inputenc}
\usepackage{amsmath,amssymb}
\usepackage{booktabs}
\usepackage{graphicx}
\usepackage{hyperref}

\title{A low-cost colorimetric sensor for lead detection}
\author{Bing Ying \and Demo Author}
\date{}
\begin{document}
\maketitle

\begin{abstract}
We report a portable colorimetric assay with a limit of detection of 0.8 ug/L.
\end{abstract}

\section{Introduction}
\label{sec:intro}

Lead contamination remains a public health problem \& requires field-ready tools (see Section\textasciitilde{}\textbackslash\{\}ref\{sec:method\}).

\section{Method}
\label{sec:method}

Samples were filtered and passed through a chelating membrane.

\begin{equation}
L = 3\sigma / S
\end{equation}

\begin{table}[htbp]
\centering
\caption{Recovery of spiked samples}
\begin{tabular}{ll}
\toprule
Matrix & Recovery (\%) \\
\midrule
Tap water & 92 \\
River water & 104 \\
\bottomrule
\end{tabular}
\end{table}

\bibliographystyle{plain}
\bibliography{refs}

\end{document}
